% ============================================================================
% Fichier  : partie3/P03_C07_famille_27000.tex
% Titre    : La Famille ISO/IEC 27000 et le SMSI
% Auteur   : MINKA MI NGUIDJOÏ Thierry Emmanuel
% ============================================================================
% Sources : ISO/IEC 27000:2018 ; 27001:2022 ; 27002, 27005. Fictif : Eséka.
% ============================================================================

\chapter{La Famille ISO/IEC 27000 et le SMSI}
\label{chap:iso27000}
\chaptermark{La famille ISO 27000}

\epigraph{%
    « Une norme ne protège rien par elle-même~:\\
    elle organise la manière dont on se protège. »%
}{Principe du cours}

\niveauChapitre{Fondamental}{La partie II.}

\begin{boxobjectifs}
\begin{itemize}
    \item Situer la famille \norme{ISO/IEC 27000} et ses principales normes.
    \item Distinguer la norme d'exigences (27001) des normes de lignes
          directrices.
    \item Comprendre ce qu'est un \sigle{SMSI} au sens de la norme.
    \item Expliquer le principe de la certification.
    \item Replacer la démarche de la mairie dans ce cadre.
\end{itemize}
\end{boxobjectifs}

\bigskip

% ============================================================================
\section*{Question directrice : pourquoi s'appuyer sur une norme plutôt que bricoler ?}
% ============================================================================

La mairie pourrait inventer sa propre méthode. S'appuyer sur
\norme{ISO/IEC 27001}, c'est hériter d'un cadre éprouvé, reconnu et auditable.
Ce chapitre présente la famille de normes et le \sigle{SMSI} qu'elle définit.

\medskip
\textbf{Plan du chapitre.}
\begin{enumerate}
    \item La famille 27000 (section~\ref{sec:c07-famille}).
    \item Exigences et lignes directrices (section~\ref{sec:c07-exigences}).
    \item Le SMSI selon la norme (section~\ref{sec:c07-smsi}).
    \item Le principe de la certification (section~\ref{sec:c07-certif}).
    \item La mairie dans le cadre (section~\ref{sec:c07-mairie}).
\end{enumerate}

\separateur

% ============================================================================
\section{La Famille 27000}
\label{sec:c07-famille}
% ============================================================================

\begin{tableaucours}{L{2.6cm}L{9.6cm}}{%
    \tableauHeader{Norme} &
    \tableauHeader{Objet}
}
ISO/IEC 27000 & Vue d'ensemble et vocabulaire \\
ISO/IEC 27001 & Exigences du SMSI (la seule certifiable) \\
ISO/IEC 27002 & Mesures de sécurité (code de bonne pratique) \\
ISO/IEC 27005 & Gestion des risques de sécurité de l'information \\
ISO/IEC 27004 & Mesure, indicateurs et évaluation \\
ISO/IEC 27006 & Exigences des organismes de certification \\
\end{tableaucours}
\vspace{-0.8em}
\begin{center}{\small\itshape Tableau~7.1. Quelques normes de la famille ISO/IEC 27000}\end{center}

La famille forme un système cohérent~: la 27001 \emph{exige}, les autres
\emph{aident}. On ne se certifie que sur la 27001~; les autres soutiennent sa
mise en \oe uvre.

% ============================================================================
\section{Exigences et Lignes Directrices}
\label{sec:c07-exigences}
% ============================================================================

\begin{boxdef}
Une norme d'\termecle{exigences} (\norme{ISO/IEC 27001}) énonce ce qui
\emph{doit} être fait pour être conforme~; ses clauses emploient « doit ». Une
norme de \termecle{lignes directrices} (27002, 27005) propose des bonnes
pratiques~; elle emploie « il convient de ».
\end{boxdef}

La distinction est capitale en audit~: un auditeur vérifie le respect des
\emph{exigences} de la 27001. Les lignes directrices éclairent le « comment »,
sans être, elles, directement certifiables.

% ============================================================================
\section{Le SMSI selon la Norme}
\label{sec:c07-smsi}
% ============================================================================

Le \sigle{SMSI} de la 27001 s'articule en sept clauses d'exigences
(clauses~4 à~10), qui suivent le cycle \sigle{PDCA}~:

\begin{tableaucours}{L{2.0cm}L{6.6cm}L{3.2cm}}{%
    \tableauHeader{Clause} &
    \tableauHeader{Exigence} &
    \tableauHeader{Étape PDCA}
}
4 & Contexte de l'organisme & Plan \\
5 & Leadership & Plan \\
6 & Planification & Plan \\
7 & Support & Do \\
8 & Fonctionnement & Do \\
9 & Évaluation des performances & Check \\
10 & Amélioration & Act \\
\end{tableaucours}
\vspace{-0.8em}
\begin{center}{\small\itshape Tableau~7.2. Les clauses d'exigences et le PDCA}\end{center}

Les Chapitres~\ref{chap:clauses46} et~\ref{chap:clauses710} détaillent ces
clauses~; le Chapitre~\ref{chap:iso27002} traite les mesures de l'Annexe~A.

% ============================================================================
\section{Le Principe de la Certification}
\label{sec:c07-certif}
% ============================================================================

\begin{boxdef}
\termecle{Certification \norme{ISO/IEC 27001}}~: attestation, délivrée par un
organisme accrédité, qu'un \sigle{SMSI} satisfait aux exigences de la norme sur
un périmètre défini. Elle se gagne par un audit et se maintient par des audits
de surveillance.
\end{boxdef}

La certification ne prouve pas l'absence d'incident~: elle atteste que
l'organisation \emph{gère} sa sécurité selon la norme. C'est une reconnaissance
de maturité, non une garantie d'invulnérabilité.

\begin{boxmetier}{Auditeur SMSI}
\textbf{Auditer l'exigence, pas l'intention.} L'auditeur ne note pas la bonne
volonté~: il vérifie, preuve à l'appui, que chaque « doit » de la norme est
satisfait. Une mairie motivée mais sans preuves échoue.
\end{boxmetier}

% ============================================================================
\section{La Mairie dans le Cadre}
\label{sec:c07-mairie}
% ============================================================================

\begin{boxprocedure}
\textbf{Constatation.} La mairie d'Eséka vise la certification \norme{ISO/IEC
27001} de son \sigle{SMSI} sur le périmètre « services aux usagers » (état
civil, taxes).

\medskip
\textbf{Analyse.} Le choix du périmètre est stratégique~: commencer par un
périmètre maîtrisable, quitte à l'étendre ensuite. Les quatre phases du fil
rouge (\textsc{Nyong}, \textsc{Kellé}, \textsc{Makak}, \textsc{Matomb})
couvrent les clauses~4 à~10.
\end{boxprocedure}

\separateur

\begin{boxresume}
\textbf{Ce qu'il faut retenir de ce chapitre~:}
\begin{itemize}
    \item La \textbf{famille 27000} forme un système~: la 27001 exige, les
          autres aident.
    \item Seule \textbf{ISO/IEC 27001} est \textbf{certifiable}~; elle emploie
          « doit ».
    \item Le \textbf{SMSI} de la norme tient en sept clauses (4 à~10) alignées
          sur le \sigle{PDCA}.
    \item La \textbf{certification} atteste la maîtrise, non l'absence
          d'incident.
    \item La mairie vise un \textbf{périmètre} défini et maîtrisable.
\end{itemize}
\end{boxresume}

\begin{boxinfo}
\textbf{Points de contrôle de l'élève.} Je sais~: citer les normes de la
famille 27000~; distinguer exigences et lignes directrices~; énoncer les sept
clauses~; expliquer la certification~; définir un périmètre.
\end{boxinfo}

\bigskip

\begin{boxexo}[title={\ding{45}\quad Exercice~7.1. La famille \niveaubase}]
Associez chaque norme à son objet~: 27001, 27002, 27005.
\end{boxexo}

\begin{boxexo}[title={\ding{45}\quad Exercice~7.2. Doit ou il convient \niveaubase}]
Laquelle de ces normes emploie « doit »~: 27001 ou 27002~? Pourquoi est-ce
important~?
\end{boxexo}

\begin{boxexo}[title={\ding{45}\quad Exercice~7.3. Les clauses \niveaubase}]
Citez les sept clauses d'exigences de la 27001 et leur étape PDCA.
\end{boxexo}

\begin{boxexo}[title={\ding{45}\quad Exercice~7.4. Ce que prouve la certification \niveaumoyen}]
« Une mairie certifiée ne peut pas être victime d'un incident. » Corrigez cette
affirmation.
\end{boxexo}

\begin{boxexo}[title={\ding{45}\quad Exercice~7.5. Rôle des lignes directrices \niveaumoyen}]
À quoi servent la 27002 et la 27005 si l'on se certifie sur la 27001~?
\end{boxexo}

\begin{boxexo}[title={\ding{45}\quad Exercice~7.6. Périmètre \niveaumoyen}]
Pourquoi la mairie commence-t-elle par un périmètre restreint~?
\end{boxexo}

\begin{boxexo}[title={\ding{45}\quad Exercice~7.7. Preuve \niveaumoyen}]
Pourquoi une mairie motivée mais sans preuves échoue-t-elle à l'audit~?
\end{boxexo}

\begin{boxexo}[title={\ding{45}\quad Exercice~7.8. PDCA et clauses \niveaumoyen}]
Reliez les clauses 9 et 10 aux étapes Check et Act, avec un exemple.
\end{boxexo}

\begin{boxexo}[title={\ding{45}\quad Exercice~7.9. Carte des normes \niveauexpert}]
Dressez une carte des normes 27000 utiles à la mairie, en indiquant pour chaque
étape du projet laquelle mobiliser.
\end{boxexo}

\begin{boxexo}[title={\ding{45}\quad Exercice~7.10. Note de périmètre \niveauexpert}]
Rédigez la note qui définit et justifie le périmètre de certification de la
mairie d'Eséka.
\end{boxexo}

\bigskip

\begin{boxinfo}
\textbf{Pour aller plus loin.}
\begin{itemize}
    \item ISO/IEC 27000~\cite{iso27000}~; ISO/IEC 27001~\cite{iso27001}.
\end{itemize}
\end{boxinfo}

\begin{boxinfo}
\textbf{Transition.} Le Chapitre~\ref{chap:clauses46} entre dans les premières
clauses~: contexte, leadership et planification.
\end{boxinfo}
